% Удиртгал

\addchaptertocentry{Удиртгал}
\pagecolor{white}

\section*{\centering \Large УДИРТГАЛ}

\subsection*{Сэдэв сонгосон үндэслэл}

Мэдээллийн технологийн хурдацтай хөгжил нь орчин үеийн иргэдийн өдөр тутмын амьдралыг ихээхэн өөрчилж байна. Цахим захиалга, онлайн үйлчилгээ нь дэлхийн өнцөг булан бүрт нэвтэрч, хүмүүсийн цаг хугацааг хэмнэж, үйлчилгээний чанарыг дээшлүүлж байна. Statista судалгааны байгууллагын 2024 оны тайланд дурдснаар дэлхийн мобайл аппликейшний зах зээлийн хэмжээ 2028 он гэхэд 756 тэрбум ам.долларт хүрнэ гэж таамаглаж байна \cite{statista_market}. Гэвч Монгол улсын орон нутгийн хотуудад, тухайлбал Өмнөговь аймгийн Даланзадгад хотод, уламжлалт үйлчилгээний хэлбэр давамгайлсаар байна.

Даланзадгад хотод өнөөдрийн байдлаар таван спорт заал тамирын цогцолбор үйл ажиллагаа явуулж байгаа бөгөөд иргэд цаг захиалахдаа дараах бэрхшээлтэй тулгардаг:

\begin{table}[h]
\centering
\caption{Одоогийн цаг захиалгын асуудлууд ба шийдэл}
\label{tab:problems}
\begin{tabularx}{\textwidth}{|p{0.4cm}|X|X|}
\hline
\textbf{\#} & \textbf{Тулгамдаж буй асуудал} & \textbf{Системийн шийдэл} \\
\hline
1 & Утсаар захиалахад шугам завгүй байх, буруу ойлголцол үүсэх & Онлайн захиалга --- цагийн хуваарь бодит цагаар харагдах \\
\hline
2 & Биечлэн очиж захиалахад цаг, зардал зарцуулагдах & Хаанаас ч гэсэн мобайл аппаар захиалах боломж \\
\hline
3 & Боломжтой цагийн мэдээлэл тодорхойгүй, ил тод бус байх & Firestore бодит цагийн шинэчлэлтээр цагийн хуваарь нэн даруй харагдах \\
\hline
4 & Захиалгаа цуцлах, өөрчлөх боломж хязгаарлагдмал байх & Апп дотроос нэг дарахад захиалга цуцлах функц \\
\hline
5 & Захиалгын баталгаажуулалт, сануулга авах боломжгүй байх & DBK-xxxxx захиалгын код + 1 цагийн өмнөх локал мэдэгдэл \\
\hline
\end{tabularx}
\end{table}

Эдгээр асуудлыг шийдвэрлэхийн тулд орчин үеийн мобайл болон \textbf{үүлэн технологи}-г ашиглан цаг захиалгын систем хөгжүүлэх нь чухал шаардлага болж байна. \keyword{Flutter} фреймворк нь нэг удаа бичиж Android болон iOS платформд нэгэн зэрэг гаргах боломж олгодог тул хөгжүүлэлтийн зардал, хугацааг мэдэгдэхүйц бууруулдаг. \keyword{Firebase} нь Google-ийн \textbf{үүлэн} бэкэнд үйлчилгээ бөгөөд бодит цагийн өгөгдөл солилцоо, найдвартай нэвтрэлт, тасралтгүй ажиллагааг хангадаг. Харин \keyword{Node.js/Express} бэкэнд нь \keyword{Docker} контейнержуулалтаар \keyword{AWS EC2} \textbf{үүлэн серверт} байршиж, \keyword{GitHub Actions} CI/CD хоолойгоор автоматаар нийтлэгддэг тул систем нь бодитоор \textbf{үүлэн орчинд} (cloud-native) ажилладаг.

%-------------------------------------------------------------------------------

\subsection*{Зорилго}

\textbf{Зорилго:} Өмнөговь аймгийн Даланзадгад хотын таван спорт заалны цаг захиалгыг цахимжуулах, хэрэглэгчдэд хялбар, хурдан, ил тод байдлаар захиалга хийх боломж олгох \keyword{Flutter} мобайл аппликейшн болон \keyword{Node.js/Express} бэкэнд бүхий гурван давхаргат системийг \keyword{үүлэн орчинд} (AWS EC2) хөгжүүлж, байршуулах.

%-------------------------------------------------------------------------------

\subsection*{Зорилт}

Уг зорилгод хүрэхийн тулд дараах зорилтуудыг дэвшүүлж байна:

\begin{enumerate}
    \item Даланзадгад хотын спорт заалнуудын цаг захиалгын одоогийн нөхцөл байдал, иргэдийн хэрэгцээг судлах, тулгамдаж буй асуудлуудыг тогтоох;
    \item Цаг захиалгын мобайл аппликейшний функциональ болон функциональ бус шаардлагын дүн шинжилгээ хийж, Use Case Diagram, Activity Diagram, Sequence Diagram боловсруулах;
    \item Flutter болон Firebase технологийг ашиглан гурван давхаргат (клиент--бэкэнд--үүлэн дэд бүтэц) системийн архитектур зохион байгуулах;
    \item Material Design 3 стандартад нийцсэн, дулаан говийн дизайн бүхий мобайл аппликейшн хэрэгжүүлэх;
    \item Node.js/Express REST API бэкэнд хөгжүүлж, Docker контейнержуулалтаар AWS EC2-д байршуулах;
    \item GitHub Actions CI/CD хоолойгоор бэкэндийн автоматжуулсан байршуулалтыг хэрэгжүүлэх;
    \item Функциональ, нэгтгэлийн болон гүйцэтгэлийн туршилт хийж, хэрэглэхэд хялбар байдлыг үнэлэх.
\end{enumerate}

%-------------------------------------------------------------------------------

\subsection*{Судлагдсан байдал}

Спорт заал болон үйлчилгээний цаг захиалгын мобайл аппликейшн хөгжүүлэлтийн чиглэлд олон улсын болон дотоодын судлаачид тодорхой хэмжээний ажил хийсэн байна.

Монгол улсын хэмжээнд \textbf{eZaal} (DAZO LLC) нь спортын байгууламжийн захиалгыг нэгтгэсэн томоохон платформ бөгөөд Google Play дээр 10,000+ татагдсан \cite{ezaal_app}. Гэхдээ зөвхөн Android платформд зориулагдсан, AI чатбот байхгүй, орон нутгийн хотуудад үйлчилгээ хүрдэггүй. Kassim et al. (2022) судалгаанд Flutter-г ашиглан sport facility booking систем хөгжүүлсэн туршлагыг дүрсэлсэн бөгөөд Firebase-тэй хослуулсан нь хөгжүүлэлтийн хугацааг 40\%-иар бууруулсан гэж тэмдэглэжээ \cite{flutter_booking}.

Монголын эрдэм шинжилгээний хэмжээнд Амарсайханы Билгүүн \cite{thesis_mobile} мобайл аппликейшн хөгжүүлэлтийн аргачлалын асуудлыг судалсан бол Т.Амартүвшин, Б.Зоригтбаатар нар \cite{mobile_mongol} мобайл апп хөгжүүлэлтийн онолын үндсийг тодорхойлсон байна. Гэхдээ тухайлан \textbf{үүлэн орчинд суурилсан}, \textbf{гурван давхаргат архитектур} бүхий спорт заалны цаг захиалгын тусгай системийн талаарх судалгаа орон нутгийн хэмжээнд хийгдээгүй байна.

Иймд энэхүү ажил нь орон нутгийн хэрэгцээнд нийцсэн, cross-platform, AI чатбот болон CI/CD автоматжуулалт бүхий орчин үеийн технологийн бүрэн стек ашигласан шинэлэг шийдлийг санал болгож байна.

%-------------------------------------------------------------------------------

\subsection*{Шинэлэг тал}

Энэхүү ажлын шинэлэг тал нь дараахид оршино:

\begin{enumerate}
    \item \textbf{Гурван давхаргат үүлэн архитектур:} Flutter клиент --- Node.js/Express бэкэнд (AWS EC2) --- Firebase/Google Cloud дэд бүтэц гэсэн тодорхой давхаргаар системийг зохион байгуулсан бөгөөд систем нь бодитоор үүлэн орчинд ажилладаг;
    \item \textbf{Docker + AWS CI/CD автоматжуулалт:} GitHub Actions ашиглан код шинэчлэгдэх бүрт Docker дүрс build хийж, AWS ECR-д push хийж, EC2-д автоматаар байршуулдаг хоолой хэрэгжүүлсэн (3--4 минут);
    \item \textbf{Давхар горимын AI чатбот:} Anthropic Claude Haiku 4.5 (Firebase Cloud Functions, үндсэн горим) болон Groq Llama 3.1-8b (Node.js бэкэнд, нөөц горим) загваруудыг нөөцлөлтийн горимоор нэгтгэж, найдвартай AI чатбот хэрэгжүүлсэн;
    \item \textbf{Байгууллагын тогтмол захиалга:} Даланзадгад хотын засгийн газар, компаниудын долоо хоногийн тогтмол цагийг \code{fixed\_bookings} цуглуулгаар удирдах тусгай функц хэрэгжүүлсэн;
    \item \textbf{Орон нутгийн хэрэгцээ:} Даланзадгад хотын таван спорт заалд тохирсон, Монгол хэлтэй, дулаан говийн дизайн бүхий мобайл аппликейшнийг Монгол улсын орон нутгийн хотод анх удаа хэрэгжүүлсэн.
\end{enumerate}

%-------------------------------------------------------------------------------

\subsection*{Судалгааны хамрах хүрээ}

\textbf{Хамрах хүрээ:} Энэхүү ажлын хүрээнд дараахи зүйлс орно:

\begin{itemize}
    \item Иргэдэд зориулсан мобайл апп (Android 8.0+ болон iOS 13.0+);
    \item Firebase дээр суурилсан email/нууц үгийн нэвтрэлт (Nodemailer 6 оронтой кодын баталгаажуулалттай, rollback механизмтай), 4 Firestore collection;
    \item Node.js/Express REST API бэкэнд (5 endpoint);
    \item Docker контейнержуулалт, AWS EC2/ECR байршуулалт (ap-southeast-1);
    \item GitHub Actions CI/CD автоматжуулалт;
    \item AI чатботын үйлчилгээ (Claude Haiku 4.5 + Groq Llama 3.1-8b);
    \item Локал мэдэгдлийн систем (\code{flutter\_local\_notifications}).
\end{itemize}

\textbf{Хамрагдахгүй зүйлс:} Онлайн төлбөрийн систем (QPay, SocialPay), GPS байршлын нарийн функц, администраторын дэлгэц.

%-------------------------------------------------------------------------------

\subsection*{Ач холбогдол}

\textbf{Практик ач холбогдол:} Даланзадгад хотын иргэдэд спорт заал захиалах үйл явцыг хялбарчилж, цаг хугацааг хэмнэнэ. Спорт заалны ажилтнуудад захиалгын менежментийг автоматжуулж, ажлын ачааллыг бууруулна. Цахим захиалга нэвтэрснээр буруу ойлголцол, шугам завгүй байхтай холбоотой асуудлуудыг арилгаж, захиалгын ил тод байдлыг хангана.

%-------------------------------------------------------------------------------

\subsection*{Ажлын бүтэц}

Энэхүү бакалаврын төгсөлтийн ажил дараахи бүтэцтэй:

\textbf{Удиртгал:} Судалгааны үндэслэл, зорилго, зорилт, судлагдсан байдал, шинэлэг тал, хамрах хүрээ, ач холбогдлыг тодорхойлсон.

\textbf{1-р бүлэг --- Онолын үндэс ба технологийн судалгаа:} Мобайл апп хөгжүүлэлтийн орчин үеийн арга хандлага, Flutter/Dart фреймворк, Firebase үйлчилгээ, Node.js/Express, Docker, AWS (EC2/ECR), GitHub Actions CI/CD, AI (Claude Haiku, Groq Llama) технологиудын онолын үндсийг судалж, ижил төстэй дотоод болон гадаадын системүүдтэй харьцуулсан.

\textbf{2-р бүлэг --- Системийн зохион байгуулалт ба хэрэгжүүлэлт:} Функциональ болон функциональ бус шаардлагын дүн шинжилгээ, Use Case Diagram, Activity Diagram, Sequence Diagram, системийн гурван давхаргат архитектур, үүлэн байршуулалтын архитектур, CI/CD хоолой, Flutter апп болон Node.js бэкэндийн дэлгэрэнгүй хэрэгжүүлэлт, Firestore өгөгдлийн загвар.

\textbf{3-р бүлэг --- Системийн туршилт ба үнэлгээ:} Функциональ туршилт (10 тохиолдол), нэгтгэлийн туршилт (4 тохиолдол), гүйцэтгэлийн хэмжилт (6 үзүүлэлт), хэрэглэхэд хялбар байдлын үнэлгээ (8 оролцогч), аюулгүй байдлын шалгалт.

\textbf{Дүгнэлт:} Ажлын нийт үр дүн, хүрсэн зорилтуудын биелэлт болон цаашдын хөгжлийн чиглэлийг дүгнэсэн.
